\section{When May a Farmer Be Accused?}
\label{sec:guarantee}

\subsection{The admissible depth is a range, not a number}
\label{sec:umax}

This layer stated a single bound, $u\le 61$\,mm, as the largest claim that never raises an alarm --- an artefact of one soil profile and one initial condition. With $V_4$ firing when
$u_i>c_{t_i}+\varepsilon$ and $c_t=d_t\cdot\mathrm{span}$, the bound at a
field with initial depletion $w_0$ irrigating no later than $d^{*}$ is
\begin{equation}
u_{\max}=\max(d^{*},w_0)\cdot\mathrm{span}+\varepsilon .
\label{eq:umax}
\end{equation}
where:
\begin{symbolist}
\item[$u_{\max}$] the largest claimed depth, in mm, that never raises an alarm;
\item[$d^{*}$] the depletion at the latest irrigated hour considered, a dimensionless fraction;
\item[$w_0$] the initial depletion of the field, in the same fraction units;
\item[$\mathrm{span}$] the width of the depletion scale, equal to $FC-WP$ in mm, defined in Sec.~\ref{sec:model};
\item[$\varepsilon$] the published tolerance for model error, equal to $5$\,mm as in Sec.~\ref{sec:filters}.
\end{symbolist}

This bound is written in depletion relative to the root zone rather than against
an absolute elevation, so unlike the absolute water-table thresholds used in current
practice it does not drift as the field surface subsides; the proposition and the
measured shift for the delta are in Supplementary~S4.

Two corrections follow, both verified by scanning the replay operator rather
than by algebra alone (derivation in Supplementary~S14): \eqref{eq:umax} uses
$\max$, not a difference --- an initially dry field admits \emph{larger}
claims than a wet one, and the difference-form expression corresponds to
no physical state --- and it bounds a \emph{single event}, not the season, so
an adversary issuing many sub-threshold claims can inflate the total without
bound. Supplementary Table~S2 gives the range that must be reported instead,
$[49.0,86.0]$\,mm across $\Theta$.

\subsection{Robust accusation under model uncertainty}
\label{sec:robust}

The auditor knows neither the farmer's soil, nor the growth stage, nor the
delivery efficiency. Let $\Theta$ be the uncertainty set and $V_\theta(L)$
the violation set under configuration $\theta$; the point-model detector
evaluated a single default $\theta_0$ and treated every violation as guilt.
We replace it by an intersection:
\begin{equation}
\begin{aligned}
V_{\mathrm{rob}}(L)&=\bigcap_{\theta\in\Theta}V_\theta(L),\\
\Theta&=\mathcal{S}\times\mathcal{K}\times\mathcal{E}\times\mathcal{W},
\end{aligned}
\label{eq:robust}
\end{equation}
where:
\begin{symbolist}
\item[$V_{\mathrm{rob}}(L)$] the robust violation set of log $L$, non-empty only if $L$ violates under every configuration;
\item[$\Theta$] the uncertainty set, the product of the four component sets below;
\item[$\mathcal{S}$] the set of soil textures, of size $5$;
\item[$\mathcal{K}$] the set of crop-coefficient regimes, of size $2$;
\item[$\mathcal{E}$] the set of irrigation efficiencies, of size $3$;
\item[$\mathcal{W}$] the set of initial states, of size $5$.
\end{symbolist}

Hence $|\Theta|=150$ and each log costs $150$ replays. The semantics are
deliberately one-sided:
\begin{equation}
a(L)=
\begin{cases}
\text{accuse}, & V_{\mathrm{rob}}(L)\neq\emptyset,\\
\text{query}, & \text{otherwise}.
\end{cases}
\label{eq:verdict}
\end{equation}
where $a(L)$ is the verdict for log $L$. A ``query'' verdict means further field
verification is required; the system does not conclude fraud on its own.

\begin{proposition}[No false accusation within $\Theta$]
\label{prop:noalarm}
If the true field configuration $\theta^{*}$ belongs to $\Theta$ and $L$ is
compliant under $\theta^{*}$, then $V_{\mathrm{rob}}(L)=\emptyset$ and the farmer
is not accused.
\end{proposition}

\begin{proof}
$V_{\mathrm{rob}}(L)\subseteq V_{\theta^{*}}(L)=\emptyset$.
\end{proof}

Proposition~\ref{prop:noalarm} is what makes the layer deployable: model error
can reduce detection power but cannot manufacture an accusation, at a price
that must be stated. Two implementation details --- the treatment of non-firing
worlds and the span of $\mathcal{W}$ --- silently break the proposition if got
wrong (Supplementary~S7).

\subsection{A proportional pump-gap threshold}

The absolute rule $G>\tau_0$ with $\tau_0=5$\,mm assumes the log is a
faithful copy of what was pumped. It is not: forgetting, rounding to the
nearest $5$\,mm and day-late entries all contribute to $G$ without wrongdoing,
so the threshold must scale with the number of recorded events:
\begin{equation}
\begin{aligned}
\tau(n)&=\delta_{\mathrm{rel}}M+z\sqrt{n}\,\sigma_{\mathrm{event}},\\
\sigma_{\mathrm{event}}^{2}&=p(1-p)\bar u^{2}+\tfrac{r^{2}}{12}
+q(1-q)(s\bar u)^{2},
\end{aligned}
\label{eq:tau}
\end{equation}
where:
\begin{symbolist}
\item[$\tau(n)$] the admissible pump gap, in mm, when the log records $n$ events;
\item[$n$] the number of irrigation events recorded in the log;
\item[$\delta_{\mathrm{rel}}$] the relative tolerance on the true pumped volume $M$, equal to $2\%$;
\item[$z$] the multiplier of the standard deviation, set to $3$;
\item[$\sigma_{\mathrm{event}}$] the standard deviation of the recording error per event, in mm;
\item[$p$] the rate of forgetting to record an event;
\item[$r$] the depth rounding step, equal to $5$\,mm;
\item[$q$] the depth-error rate, and $s$ its relative amplitude;
\item[$\bar u$] a typical event depth, equal to $30$\,mm.
\end{symbolist}

The two rules are not two cases of one formula, and \eqref{eq:tau} should not be
read as a generalisation of $\tau_0$. As $n\to0$ and
$\sigma_{\mathrm{event}}\to0$ the expression above tends to
$\delta_{\mathrm{rel}}M$, that is $2\%$ of the metered volume, not to $5$~mm; the two
thresholds coincide only if $M$ happens to equal $250$~mm. They differ in kind:
$\tau_0=5$~mm is an absolute constant that is valid only if the log is assumed to
be an exact copy of the pumped volume, whereas $\tau(n)$ scales with both $M$ and
the event count $n$. In the experiments every verdict uses $\tau(n)$, and
$\tau_0$ is kept only as a reference point that exposes the cost of assuming a
faithful copy.

\begin{remark}[A filter that cannot fire]
\label{sec:v2reach}
The third physical filter, flagging depletion below the wilting point, was doubly inoperative: its threshold ($1.02$) lay outside the clipped range of the reported indicator ($0$ firings in $281$ logs), and once repaired it still cannot activate --- across $150$ configurations with an empty log the unclipped depletion peaks at $0.9987$, and a dry-season delta rice field does not reach the wilting point (Supplementary Table~S3, derivation in
S13). The layer has two operative physical filters and one inert safety device.
\end{remark}
